% GENERATED FILE. Edit canonical lessons, then run npm run generate:books.
% canonical-source-hash: fnv1a64:a363165fe1eeeef0
% canonical-lessons: BN-C265-jabo, BN-C265-shesh, BN-W265-letter-murdhanya-sha, BN-R265-what-will-happen

\chapter{What Will Happen, and the End}
\label{ch:bn-what-will-happen-and-the-end}
\hlchaptermodality{265}

\begin{chapteropening}
\textbf{By the end of this chapter:} \emph{I can say what will happen and what I will do (jābô, jābe, jāben, hôbe), say what I won't do, and read and write \bn{ষ} in \bn{শেষ}.}

Complete BN-R265-what-will-happen: say what will happen, what you will do and what you won't.
\end{chapteropening}

\section[\texorpdfstring{\bn{আমি} \bn{যাব} (āmi jābô)}{আমি যাব (āmi jābô)}]{\bn{আমি} \bn{যাব} (āmi jābô) — the future, person by person}

\label{lesson:BN-C265-jabo}



\begin{quote}
Say \textquotedblleft{}Yesterday I was at home\textquotedblright{}. (\emph{kāl āmi bāṛite chhilām}.) Then the goodbye from chapter seven: \textbf{\bn{কাল} \bn{দেখা} \bn{হবে}}, \textquotedblleft{}see you tomorrow\textquotedblright{}. Its last word, \textbf{\bn{হবে}}, is a future: \textquotedblleft{}it will happen\textquotedblright{}.
\end{quote}

\subsection*{You'll want to know: the future, person by person}
\begin{quote}
\textbf{\bn{কাল} \bn{আমি} \bn{স্টেশনে} \bn{যাব।}}
\end{quote}

\begin{quote}
\emph{kāl āmi sṭeshône jābô.} — \textquotedblleft{}Tomorrow I will go to the station.\textquotedblright{}
\end{quote}

With a future, \textbf{\bn{কাল}} is tomorrow. The future puts \textbf{-\bn{ব}-} on the stem, then the person:

\begin{itemize}
\raggedright
  \item \textbf{\bn{আমি}} / \textbf{\bn{আমরা}} \textbf{\bn{যাব}} (\emph{jābô}) — I will go, we will go
  \item \textbf{\bn{তুমি}} / \textbf{\bn{সে}} \textbf{\bn{যাবে}} (\emph{jābe}) — you will go; he or she will go
  \item \textbf{\bn{আপনি}} / \textbf{\bn{তিনি}} \textbf{\bn{যাবেন}} (\emph{jāben}) — with respect
  \item \textbf{\bn{তোমরা}} / \textbf{\bn{তারা}} \textbf{\bn{যাবে}} — you, to several; they
\end{itemize}

\textbf{\bn{তুমি}} and \textbf{\bn{সে}} share \textbf{-\bn{বে}} here; the past kept them apart (\textbf{\bn{গেলে}}, \textbf{\bn{গেল}}). The stem is the plain one, not the past's: \textbf{\bn{করব}} (\emph{kôrbô}), \textbf{\bn{দেখব}} (\emph{dekhbô}), \textbf{\bn{খাব}} (\emph{khābô}), \textbf{\bn{আসব}} (\emph{āsbô}), \textbf{\bn{থাকব}} (\emph{thākbô}); \textbf{\bn{দেওয়া}} $\to$ \textbf{\bn{দেব}} (\emph{debô}), \textbf{\bn{নেওয়া}} $\to$ \textbf{\bn{নেব}} (\emph{nebô}); \textbf{\bn{হওয়া}} $\to$ \textbf{\bn{হবে}} (\emph{hôbe}), \textquotedblleft{}will be, will happen\textquotedblright{}. You have said one \textbf{-\bn{বেন}} for a long time: \textbf{\bn{মাফ} \bn{করবেন}}, \textquotedblleft{}excuse me\textquotedblright{}.

\subsection*{You'll want to know: will not}
\textquotedblleft{}Will not\textquotedblright{} puts \textbf{\bn{না}} after the verb: \textbf{\bn{আমি} \bn{যাব} \bn{না}} (\emph{āmi jābô nā}), \textquotedblleft{}I won't go\textquotedblright{}.

\begin{quote}
\textbf{\bn{কাল} \bn{আমি} \bn{বাজারে} \bn{যাব} \bn{না}, \bn{বাড়িতে} \bn{থাকব।}}
\end{quote}

\begin{quote}
\emph{kāl āmi bājāre jābô nā, bāṛite thākbô.} — \textquotedblleft{}Tomorrow I won't go to the market; I'll stay at home.\textquotedblright{}
\end{quote}

\subsection*{Your turn}
Say these aloud:

\begin{itemize}
\raggedright
  \item \textquotedblleft{}Tomorrow I will go to the station\textquotedblright{} — \emph{kāl āmi sṭeshône jābô}
  \item \textquotedblleft{}Will you come?\textquotedblright{}, to a friend — \emph{tumi āsbe?}
  \item \textquotedblleft{}What will happen?\textquotedblright{} — \emph{kī hôbe?}
\end{itemize}

\begin{tcolorbox}[breakable,colback=teal!4,colframe=teal!35!black,title={Before you move on}]
\textquotedblleft{}I will go\textquotedblright{}? (\textbf{\bn{যাব}}.) \textquotedblleft{}She will come\textquotedblright{}? (\textbf{\bn{আসবে}}.) \textquotedblleft{}You will give\textquotedblright{}, with respect? (\textbf{\bn{দেবেন}}.) \textquotedblleft{}I won't go\textquotedblright{}? (\textbf{\bn{যাব} \bn{না}}.)
\end{tcolorbox}
\section[\texorpdfstring{\bn{শেষ} (shesh)}{শেষ (shesh)}]{\bn{শেষ} (shesh) — the end; last}

\label{lesson:BN-C265-shesh}



\begin{quote}
Before the new one: say \textquotedblleft{}I will go\textquotedblright{} and \textquotedblleft{}I won't go\textquotedblright{}. (\emph{jābô}, \emph{jābô nā}.)
\end{quote}

\subsection*{You'll want to know: \bn{শেষ}}
\textbf{\bn{শেষ}} — \emph{shesh} — \textquotedblleft{}the end; last\textquotedblright{}.

\begin{quote}
\textbf{\bn{কাল} \bn{আমার} \bn{কাজ} \bn{শেষ} \bn{হবে।}}
\end{quote}

\begin{quote}
\emph{kāl āmār kāj shesh hôbe.} — \textquotedblleft{}My work will be finished tomorrow.\textquotedblright{}
\end{quote}

\textbf{\bn{শেষ} \bn{হওয়া}} is \textquotedblleft{}to come to an end\textquotedblright{}, and \textbf{\bn{শেষ} \bn{করা}} \textquotedblleft{}to finish\textquotedblright{}: \textbf{\bn{আমি} \bn{বইটা} \bn{শেষ} \bn{করলাম}} (\emph{āmi bôiṭā shesh kôrlām}), \textquotedblleft{}I finished the book\textquotedblright{}. Before a noun, \textbf{\bn{শেষ}} is \textquotedblleft{}last\textquotedblright{}:

\begin{quote}
\textbf{\bn{শেষ} \bn{ট্রেন} \bn{রাতে} \bn{ছাড়বে।}}
\end{quote}

\begin{quote}
\emph{shesh ṭren rāte chhāṛbe.} — \textquotedblleft{}The last train will leave at night.\textquotedblright{}
\end{quote}

\begin{cousinweb}
\textbf{\bn{শেষ}} (\emph{shesh}) is a tatsama: Sanskrit \emph{śeṣa}, \textquotedblleft{}what remains\textquotedblright{}, which CDIAL (12611) and Monier-Williams both give. Monier-Williams also lists \textquotedblleft{}end, conclusion\textquotedblright{} and \textquotedblleft{}last\textquotedblright{} among its senses, so the Bengali uses were already Sanskrit's.
\end{cousinweb}

\subsection*{Your turn}
Say these aloud:

\begin{itemize}
\raggedright
  \item \emph{shesh}
  \item \textquotedblleft{}My work will be finished tomorrow\textquotedblright{} — \emph{kāl āmār kāj shesh hôbe}
  \item \textquotedblleft{}The last train will leave at night\textquotedblright{} — \emph{shesh ṭren rāte chhāṛbe}
\end{itemize}

\begin{tcolorbox}[breakable,colback=teal!4,colframe=teal!35!black,title={Before you move on}]
What does \textbf{\bn{শেষ}} mean? (\textquotedblleft{}The end\textquotedblright{}, and before a noun, \textquotedblleft{}last\textquotedblright{}.) \textquotedblleft{}To finish\textquotedblright{}? (\textbf{\bn{শেষ} \bn{করা}}.)
\end{tcolorbox}
\section[\texorpdfstring{\bn{ষ} (sh)}{ষ (sh)}]{\bn{ষ} (sh) — the letter at the end of \bn{শেষ}}

\label{lesson:BN-W265-letter-murdhanya-sha}



\begin{quote}
Before the new one: say \textbf{\bn{শেষ}}, \textquotedblleft{}the end\textquotedblright{}. Its first letter is \textbf{\bn{শ}}, which you write; its last is new.
\end{quote}

\subsection*{Script you'll notice: \bn{ষ}}
\textbf{\bn{ষ}} — \emph{sh}. Bengali says \textbf{\bn{শ}}, \textbf{\bn{ষ}} and \textbf{\bn{স}} alike, as \emph{sh} in most words, so the letter is a matter of spelling: \textbf{\bn{শেষ}} begins with \textbf{\bn{শ}} and ends with \textbf{\bn{ষ}}. Most of the words that carry \textbf{\bn{ষ}} come from Sanskrit.

In the print of this book, \textbf{\bn{ষ}} has a headline and an upright stem on the right, the full height of the letter, as \textbf{\bn{ব}} and \textbf{\bn{ঘ}} do. A long straight line lies on a slant between the left end of the headline, close under it, and the middle of the stem. A shorter slanting line joins the middle of that line to a pointed tip at the lower left, and a curve joins that tip to the foot of the stem.

\subsection*{Writing: \bn{ষ} — copy the printed shape}
Put your pen beside \textbf{\bn{ষ}} and copy the shape you can see, slowly and large. Check the copy against the print: the headline, the stem, the long slanting line, the short one to the tip, and the curve to the foot.

\begin{quote}
This book does not yet tell you where to start \bn{ষ} or which way the pen goes. That is not written down here until it can be written down with a source. Copying the printed shape needs no such source, and it is how the shape becomes yours in the meantime.
\end{quote}

\subsection*{Your turn}
\begin{itemize}
\raggedright
  \item \emph{Find:} \bn{ষ} at the end of \textbf{\bn{শেষ}}
  \item \emph{Write it:} \bn{শ}, then \bn{ষ}, from the print, naming each
  \item \emph{Write it:} \textbf{\bn{শেষ}}
  \item \emph{Write it:} \bn{ঘ}, from memory
\end{itemize}

\begin{tcolorbox}[breakable,colback=teal!4,colframe=teal!35!black,title={Before you move on}]
Which letter ends \textbf{\bn{শেষ}}? (\textbf{\bn{ষ}}.) How is it said? (\emph{sh}, as \textbf{\bn{শ}} is.) What lies across its middle? (A long slanting line, meeting the stem.)
\end{tcolorbox}
\section[(review)]{Review — what will happen}

\label{lesson:BN-R265-what-will-happen}



\begin{quote}
Say \textquotedblleft{}Yesterday I went to the market\textquotedblright{}, then \textquotedblleft{}Tomorrow I will go to the station\textquotedblright{}. (\emph{kāl āmi bājāre gelām}; \emph{kāl āmi sṭeshône jābô}.)
\end{quote}

\subsection*{Your turn}
\begin{itemize}
\raggedright
  \item \emph{Say it:} Tomorrow we will eat fish. Mira will come too.
  \item \emph{Say it:} I won't go to the market; I'll stay at home.
  \item \emph{Say it:} When will the last train leave?
  \item \emph{Say it:} Yesterday I was tired. Tomorrow I'll rest.
  \item \emph{Say it:} My work will end in the evening.
  \item \emph{Write it:} \textbf{\bn{শেষ}}, then \bn{ঘ}
\end{itemize}

\begin{tcolorbox}[breakable,colback=teal!4,colframe=teal!35!black,title={Before you move on}]
Which two people share \textbf{-\bn{বে}}? (\textbf{\bn{তুমি}} and \textbf{\bn{সে}}.) \textquotedblleft{}It will happen\textquotedblright{}? (\textbf{\bn{হবে}}.) \textquotedblleft{}Will not\textquotedblright{}? (\textbf{\bn{না}} after the verb.) Which letter ends \textbf{\bn{শেষ}}? (\textbf{\bn{ষ}}.)
\end{tcolorbox}
